\section{Background}
\subsection{Linear Temporal Logic over finite traces}

Linear Temporal Logic (LTL) \cite{LTL} is a formal language that extends traditional propositional logic with modal operators, allowing the specification of rules that must hold \textit{through time}. Here, we use \ltl over finite traces (\ltlf) \cite{LTLf}, a popular variant of LTL which models finite, but length-unbounded, traces of executions, making it suitable for finite-horizon problems.
%%
Given a finite set $\Sigma$ of atomic propositions, the set of \ltlf formulas $\varphi$ is inductively defined as follows:
%\begin{equation}
\[
    \varphi ::= \top \mid\bot\mid \sigma\mid\lnot\varphi\mid\varphi_1\land\varphi_2\mid X\varphi\mid\varphi_1 U\varphi_2,
\]
%\end{equation}
where $\sigma \in \Sigma$. We use $\top$ and $\bot$ to denote $True$ and $False$, respectively. $X$ (Strong Next) and $U$ (Until) are temporal operators. Other temporal operators are $N$ (Weak Next) and $R$ (Release), defined as $N \varphi \equiv \neg X\neg\varphi$ and $\varphi_1 R \varphi_2 \equiv \neg(\neg\varphi_1 U \neg\varphi_2 )$; $G$ (globally) $G\varphi \equiv \bot R\varphi$; and $F$ (eventually) $F \varphi \equiv \top U \varphi$.


A trace $\boldsymbol{\sigma} = (\sigma_{1}, \sigma_{2}, \dots, \sigma_{T})$ is a sequence of propositional assignments to the propositions in $\Sigma$, where $\sigma_{t} \in 2^\Sigma$ is the set of all and only propositions that are true at instant $t$. Additionally, $|\boldsymbol{\sigma}| = T$ denotes the length of the trace. Since every trace is finite, $|\boldsymbol{\sigma}| < \infty$.
If the propositional symbols in $\Sigma$ are all \textit{mutually exclusive}, i.e., the domain produces exactly one symbol true at each step, then we have $\sigma_{t} \in \Sigma$.
%
By $\boldsymbol{\sigma} \vDash \varphi$ we denote that the trace $\boldsymbol{\sigma}$ satisfies the \ltlf formula $\varphi$. 
We refer the reader to \cite{LTLf} for a formal description of the \ltlf semantics. 

Any \ltlf formula $\varphi$ can be translated into a Deterministic Finite Automaton (DFA) \cite{LTLf} $A_\varphi = (2^\Sigma, Q, q_0, \delta, F)$, where $2^\Sigma$ is the automaton alphabet, $Q$ is the finite set of states, $q_0 \in Q$ is the initial state, $\delta: Q \times 2^\Sigma \rightarrow Q$ is the transition function, and $F \subseteq Q$ is the set of final states. 
Additionally, we recursively define the extended transition function over traces $\delta^*: Q \times (2^\Sigma)^* \rightarrow Q$ as:

\begin{equation}
%\[
\begin{array}{l}
    \delta^*(q,\epsilon) = q \\
    \delta^*(q, \sigma+\boldsymbol{x}) = \delta^*(\delta(q,\sigma) , \boldsymbol{x}),
\end{array}
%\]
\end{equation}
where $\sigma \in 2^\Sigma$ is a symbol and $\boldsymbol{x} \in (2^\Sigma)^*$ is a trace.  
The automaton accepts the trace $\boldsymbol{\sigma}$ if $\delta^*(q_0, \boldsymbol{\sigma}) \in F$, and in that case we say that $\boldsymbol{\sigma}$ belongs to the language of the automaton, denoted as $L(A_\varphi)$.  
We have that $\varphi$ and $A_\varphi$ are equivalent because, for any trace $\boldsymbol{\sigma} \in (2^\Sigma)^*$:
\begin{equation}
%\[
\boldsymbol{\sigma}\in L(A_\varphi) \iff\boldsymbol{\sigma}\vDash\varphi.
%\]
\end{equation}

\subsection{Offline Reinforcement Learning}

In Reinforcement Learning (RL), the interaction between an agent and the environment is commonly modeled as a Markov Decision Process (MDP) \cite{sutton1998reinforcement}, defined by the tuple $(S,A,t,r,\gamma)$ where $S$ is the set of states, $A$ the set of actions, $t : S \times A \times S \rightarrow [0,1]$ the transition function, $r : S \times A \rightarrow R$ the reward function, and $\gamma \in [0,1]$ the discount factor. A policy $\pi : S \rightarrow A$ maps states to actions, while the value function $V^\pi(s)$ denotes the expected discounted return obtained by following $\pi$ from state $s$. The goal of the RL agent is to learn the optimal policy $\pi^*$ that provides the maximum value. 

\begin{comment}
Standard MDPs assume
Markovian transitions and rewards, i.e., dependence only on
the current state. However, many real-world problems violate
this assumption [Littman et al., 2017]. In a non-Markovian
decision process, the reward function r : (S × A)∗ → R,
the transition function t : (S × A)∗ × S → [0,1], or both
may depend on the interaction history. In this work, we focus
on Non-Markovian Reward Decision Processes (NMRDPs)
[Bacchus et al., 1996], where non-Markovianity arises solely
from the reward function. Learning optimal policies in NM
RDPsis challenging, as rewards may depend on the entire se
quence of past state–action pairs, making standard RL meth
ods inapplicable. To address this issue, many approaches con
struct an augmented Markovian state representation by moni
toring the task through a labeling function L : S → P, which
maps environment states to propositional symbols. The re
sulting labeled traces are used to track task progress, typi
cally via LTL formula progression [Toro Icarte et al., 2018;
Vaezipoor et al., 2021] or equivalent automata-based repre
sentations such as Moore machines [Camacho et al., 2019;
De Giacomo et al., 2019].
\end{comment}

Traditional RL methods assume direct interaction with the environment during training. The agent collects experience by executing actions, observing state transitions and rewards, and updating its policy accordingly. This online interaction paradigm enables iterative improvement but may be costly, unsafe or even unfeasible in real-world applications. This motivates the study of offline RL \cite{levine2020offline}, where the agent is given a fixed dataset $\D$ of trajectories generated by some (unknown) behavior policy, and must learn a policy without additional environment interaction. Many offline RL algorithms incorporate mechanisms to constrain learned policies to remain close to the support of the data distribution. %This is necessary to mitigate distributional shift, since selecting actions that are poorly represented in the dataset may lead to inaccurate value estimates and degraded policy performance.
%\textcolor{blue}{Ele: fino a qui basta, non abbiamo spazio per parlare di model free e model based}
%In general, offline RL is addressed by two different families of approaches: model-free RL and model-based RL.
%\textcolor{blue}{TODO: Complete it and put references.}

\subsection{Reinforcement Learning as Sequence Modeling}

Some works \cite{trajectory_transformers,decision_transformers} propose an alternative perspective on RL, viewing it as a sequence modeling problem rather than a value estimation. In \cite{trajectory_transformers} a trajectory is treated as a sequence of tokens and an autoregressive model is used to learn the next token in the sequence.
%
Suppose we have $N$-dimensional states and $M$-dimensional actions. A trajectory $\boldsymbol{\tau}$ of horizon $T$ can be represented as
\begin{equation}
\label{eq:trace_form}
%\[
\boldsymbol{\tau} = (s_1^1, s_1^2, ..., s_1^N, a_1^1, ..., a_1^M, r_1, s_2^1, ..., r_T)
%\]
\end{equation}
%
\noindent
Subscripts on all tokens denote timestep and superscripts on states and actions denote dimension (i.e., $s_t^i$ is the $i^\text{th}$ dimension of the state at time $t$). An autoregressive model is capable of estimating the probability of the $t^\text{th}$ token $x_{t}$ given the set of all past tokens $\boldsymbol{x_{<t}}$:
\begin{equation}
\label{eq:next_act_pred}
P(x_{t} \mid x_{1}, x_{2}, \dots, x_{t-1}) = P(x_{t} \mid \boldsymbol{x_{<t}}).
\end{equation}
%
\noindent
Instead of learning value functions or policies, one can learn the joint distribution
\begin{equation}
P(\boldsymbol{x}) = \prod_{i=1}^{|\tau|} P(x_{i} \mid \boldsymbol{x_{<i}}).
\end{equation}

At each generation step, a token must be \textit{sampled} from the next activity probability. A common way of selecting the next token to feed into the autoregressor is to \textit{greedily} choose it by maximizing the next token probability at each step, as follows:
%\[
\begin{equation}
{x}_{t} = \argmax_x P(x_{t} \mid \boldsymbol{x}_{<t}), \quad \forall t < |\tau|.
\end{equation}
%\]
\noindent
In general, this greedy search strategy is sub-optimal, because it may not produce the trace that maximizes the joint probability distribution. Other sub-optimal sampling strategies commonly used for this task include Beam Search, Random Sampling, and Temperature Sampling~\cite{ramamaneiro24}.

Trajectory Transformer (TT) models are transformer-based autoregressive architectures well suited for addressing RL as sequence modeling. In the offline setting, such models can be trained via maximum likelihood using teacher forcing, by maximizing the log-probability of trajectories in the dataset. At inference time, planning can be performed by sampling or beam search, optionally biasing candidate trajectories according to cumulative reward or reward-to-go estimates. This paradigm unifies dynamics modeling, policy learning, and behavior regularization within a single sequence model.

Note that this approach focuses exclusively on maximizing return and modeling the data distribution, but they may violate safety or liveness constraints that naturally emerge from (or we want to impose to) the domain. This motivates our neurosymbolic framework guiding the learned trajectory distribution toward constraint satisfaction.